\documentclass[11pt]{article}
\usepackage[letterpaper,margin=0.68in,headheight=15pt]{geometry}
\usepackage[T1]{fontenc}
\usepackage{lmodern}
\usepackage{amsmath,amssymb}
\usepackage{array,booktabs,enumitem,multicol}
\usepackage{xcolor}
\usepackage{tikz}
\usepackage{fancyhdr}
\usepackage{microtype}
\setlength{\parindent}{0pt}
\setlength{\parskip}{5pt}
\definecolor{olive}{HTML}{555B35}
\definecolor{gold}{HTML}{B18A22}
\definecolor{soft}{HTML}{F4F1DE}
\definecolor{bluegray}{HTML}{E8EEF2}
\pagestyle{fancy}
\fancyhf{}
\fancyhead[L]{\textcolor{olive}{\small MATH\& 141: Precalculus I}}
\fancyhead[R]{\textcolor{gold}{\small Fall 2026}}
\fancyfoot[C]{\thepage}
\renewcommand{\headrulewidth}{0.45pt}
\renewcommand{\footrulewidth}{0pt}
\setlist[enumerate]{leftmargin=*,itemsep=0.08in,topsep=0.05in}
\newcommand{\assessmenttitle}[2]{%
  {\Large\bfseries Module #1: #2}\par\vspace{3pt}
}
\newcommand{\studentline}{%
\noindent\textbf{Name:}\hspace{2.6in}\textbf{Date:}\hspace{1.15in}\par\vspace{7pt}}
\newcommand{\directions}[1]{\textit{#1}\par\vspace{4pt}}
\newcommand{\answerblank}[1]{\vspace{#1}}
\newcommand{\blankgrid}[4]{%
\begin{center}
\begin{tikzpicture}[x=0.75cm,y=0.75cm]
  \draw[step=1,gray!28,very thin] (#1,#3) grid (#2,#4);
  \draw[->,thick] (#1,0)--(#2,0) node[right] {$x$};
  \draw[->,thick] (0,#3)--(0,#4) node[above] {$y$};
  \foreach \x in {#1,...,#2}{\ifnum\x=0\else\draw (\x,0.10)--(\x,-0.10) node[below=2pt,font=\scriptsize]{\x};\fi}
  \foreach \y in {#3,...,#4}{\ifnum\y=0\else\draw (0.10,\y)--(-0.10,\y) node[left=2pt,font=\scriptsize]{\y};\fi}
\end{tikzpicture}
\end{center}}
\newcommand{\smallgrid}[4]{%
\begin{tikzpicture}[x=0.50cm,y=0.50cm]
  \draw[step=1,gray!28,very thin] (#1,#3) grid (#2,#4);
  \draw[->,thick] (#1,0)--(#2,0) node[right] {$x$};
  \draw[->,thick] (0,#3)--(0,#4) node[above] {$y$};
\end{tikzpicture}}
\begin{document}

% MODULE 9 PREQUIZ
\assessmenttitle{9}{Systems of Linear Equations -- Prequiz}
\studentline
\directions{Four short questions. Interpret a solution as a common point satisfying every equation.}
\begin{enumerate}
\item Solve the system $x+y=8$ and $x-y=2$.
\answerblank{0.52in}
\item Classify the system $2x+y=4$ and $4x+2y=10$ as having one solution, no solution, or infinitely many solutions. Explain briefly.
\answerblank{0.62in}
\item Solve the triangular system
\[
z=2,\qquad y-z=1,\qquad x+y+z=9.
\]
\answerblank{0.72in}
\item What is the geometric meaning of a unique solution to two linear equations in two variables?
\answerblank{0.48in}
\end{enumerate}
\newpage

% MODULE 9 CHECKIN PAGE 1
\assessmenttitle{9}{Systems of Linear Equations -- Weekly Check-In}
\studentline
\directions{Three questions. Interpret a solution as a common point satisfying every equation.}
\textbf{1. Solve and interpret a two-variable system.}

Solve
\[
\begin{cases}
2x+y=7,\\
x-y=2.
\end{cases}
\]
\begin{enumerate}[label=(\alph*),leftmargin=0.35in]
\item Solve algebraically and state the ordered-pair solution.
\answerblank{0.85in}
\item Verify the ordered pair in both original equations.
\answerblank{0.55in}
\item What does this solution represent geometrically?
\answerblank{0.42in}
\end{enumerate}

\textbf{2. Classifying systems.}

Classify each system as having one solution, no solution, or infinitely many solutions. Explain using the relationship between the equations.

\textbf{(a)} $2x+3y=6$ and $4x+6y=12$.\hfill
\textbf{(b)} $2x+3y=6$ and $4x+6y=10$.
\answerblank{1.05in}
\newpage

% MODULE 9 CHECKIN PAGE 2
\assessmenttitle{9}{Systems of Linear Equations -- Weekly Check-In, continued}
\studentline
\textbf{3. Back-substitution in three variables.}

Solve
\[
\begin{cases}
z=2,\\
y-3z=-1,\\
2x+y+z=15.
\end{cases}
\]
Show the order in which you solve for the variables. Then check your ordered triple in all three equations.
\answerblank{4.35in}

\end{document}
